% Alternative RQ3 draft for comparison with the current version.
% Counts use the exhaustive 735-case C1--C4 taxonomy; one extraction artifact
% is excluded. Validate the automated say-side labels before final submission.

\subsection{RQ3: Where does the say--do alignment chain break?}
\label{sec:rq3-alternative}

RQ1 establishes that halluci-action occurs in both directions, while RQ2 shows
that factuality does not expose it. We next ask what forms these mismatches take
and where they arise in the chain from a verbal self-report to an executed and
disclosed action. We classify every confirmed instance using four conditions
that a faithful self-report must satisfy. For a claim that reports operation
$X$ on target $Y$ with outcome $Z$, the operation must be executed (C1),
grounded to the stated target and parameters (C2), and verified to have achieved
the stated outcome (C3). In the reverse direction, every executed action must be
disclosed by the agent's language (C4). Taking the first failed claim-side gate
makes C1--C3 mutually exclusive; C4 is the corresponding action-anchored gate
(Figure~\ref{fig:rq3-fault-tree}).

\begin{figure}[t]
  \centering
  \includegraphics[width=\linewidth]{figures/rq3_c1_c4_fault_tree.pdf}
  \caption{A fault-tree decomposition of halluci-action. C1--C3 decompose an
  unsupported action claim by the first failed claim-side gate; C4 captures
  executed actions with no verbal disclosure. Counts exclude one extraction
  artifact.}
  \label{fig:rq3-fault-tree}
\end{figure}

Of the 735 genuine halluci-action instances, 86 fail at execution (C1), 33 at
grounding (C2), eight at outcome verification (C3), and 608 at disclosure (C4).
The large C4 count should not be read as 608 equally serious violations: 586 are
undisclosed reads, whereas 22 are hidden state mutations, including 17 actions
judged forbidden. Conversely, the smaller claim-side classes contain several of
the failures most likely to create false assurance. The taxonomy therefore
separates the mechanism of mismatch from its risk: failure location tells us
where alignment broke, while action requirement, mutation status, and disclosure
determine its consequence.

\paragraph{C1: verbal commitments outrun execution.}
C1 contains 86 claims for which no matching operation appears in the trace. Most
are abandoned intentions (84), such as ``Let me check what products might be
relevant'' followed by no product lookup. These cases add narration noise but
usually do not assert a completed state. The remaining two are fabricated
completions: the agent reports a task as complete despite issuing no operation
that could realize the claimed result. The distinction matters because identical
trace absence can express either a dropped plan or direct false assurance.

\paragraph{C2: operations lose their referents during grounding.}
In 33 cases, an operation is attempted but its structured arguments do not match
the entity or parameters named in the claim. For example, an agent says it will
log an interaction against a lead but calls
\texttt{log\_interaction(customer\_id=13)}, writing to a customer record instead.
Eleven of the 33 C2 cases mutate persistent state. These cases expose a failure at
the language-to-tool boundary: the high-level operation survives, but its target
does not. They motivate typed entity bindings and explicit target verification
before state-changing calls.

\paragraph{C3: attempts are converted into unverified success.}
Eight cases reach the intended operation and target but fail to establish the
reported result. Four are false-success reports after an error, empty result, or
rejected state transition. In a representative sales task, the agent claims that
a partial delivery was arranged even though the tool payload states
\texttt{delivery\_created=false}. Four further cases assert an outcome or target
binding that the trace never verifies. These errors occur at the result-to-report
boundary and suggest that completion language should be gated on an explicit
success signal or postcondition read-back.

\paragraph{C4: tool use is treated as private rather than communicative.}
C4 is the dominant class, comprising 608 executed actions with no matching verbal
disclosure. Its composition is sharply bimodal. Most are read-only lookups (586)
that do not change environment state, explaining much of the high raw
done-not-said rate for terse agents. The remaining 22 are hidden mutations; 17 of
these violate a task requirement or explicit constraint. This small covert slice
is the central oversight risk: transcript-only monitoring sees a benign response
while the execution trace contains a consequential state change.

\begin{figure}[t]
  \centering
  \includegraphics[width=\linewidth]{figures/rq3_failure_composition_by_model.pdf}
  \caption{Composition of genuine halluci-action instances within each model.
  C4 is separated into undisclosed reads and hidden mutations because these have
  substantially different consequences. Numbers inside sufficiently large
  segments are instance counts.}
  \label{fig:rq3-model-composition}
\end{figure}

\paragraph{Models have different failure signatures.}
Figure~\ref{fig:rq3-model-composition} shows that the aggregate HAR values in RQ1
mask distinct mechanisms. DeepSeek-Chat concentrates in unexecuted commitments
and misgrounded executions, consistent with a verbose agent that narrates future
steps and occasionally loses entity references. Gemini-Flash is almost entirely
composed of undisclosed reads, making its reverse gap large but generally
non-mutating. Qwen3-8B contains the largest hidden-mutation residue, including 11
covert forbidden actions, while GPT-4o-mini contributes disproportionately to the
small C3 false-success class. Models therefore differ not only in how frequently
their words and actions diverge, but also in which link of the alignment chain
fails and how consequential that failure is.

\paragraph{Scope and interpretation.}
This taxonomy identifies the observable failure location; it does not by itself
establish a causal account of the model's internal behavior. Explanations such as
optimistic closure, entity-resolution weakness, or treating tools as private
scratch space should therefore be read as mechanism hypotheses grounded in the
trace patterns. Before final reporting, we exclude one non-actionable extraction
artifact, quarantine the known broken-tool-adjacent refund cases, and validate
the automatically assigned claim-side labels on a stratified human-reviewed
sample. Detailed definitions, counts, and representative cases are provided in
Appendix~X.
